\chapter{文件系统}

\section{文件基本概念}
\concept{文件系统的组成}
\begin{points}
  \item 文件系统是操作系统中与管理文件有关的软件和数据的集合。
  \item 组成包括：被管理的文件集合、管理文件所需的数据结构、相应管理软件。
  \item 管理软件负责外存空间分配回收、目录管理、文件读写、共享与保护等。
\end{points}
\plain{文件是用户看到的逻辑外存单位；用户不用管扇区，只管名字、属性和读写操作。}
\exam{常考文件属性、文件操作、open 为什么能减少目录搜索次数。}


\concept{文件定义}
\begin{points}
  \item 文件是记录在外部存储上，具有名字，在逻辑上有完整意义的信息项序列。
  \item 信息项是构成文件内容的基本单位，可以是字节或记录。
  \item 从用户角度看，文件是逻辑外存的最小分配单元。
  \item 文件内容的意义由文件创建者和使用者解释。
\end{points}
\plain{文件是用户看到的逻辑外存单位；用户不用管扇区，只管名字、属性和读写操作。}
\exam{常考文件属性、文件操作、open 为什么能减少目录搜索次数。}


\concept{文件属性}
\begin{points}
  \item 名称：便于用户识别的符号文件名。
  \item 标识符：系统内部唯一标识，常为数字形式。
  \item 类型：普通文件、目录文件、设备文件等。
  \item 位置：文件在设备上的位置指针。
  \item 大小：当前大小和最大允许大小。
  \item 保护：访问控制信息，规定谁能读、写、执行。
  \item 时间、日期、用户标识：创建、修改、访问等元数据。
\end{points}
\plain{文件是用户看到的逻辑外存单位；用户不用管扇区，只管名字、属性和读写操作。}
\exam{常考文件属性、文件操作、open 为什么能减少目录搜索次数。}


\concept{文件操作}
\begin{points}
  \item Create：创建文件，为文件分配目录项和必要空间。
  \item Write：根据写指针向文件写入数据。
  \item Read：根据读指针从文件读取数据。
  \item Seek：移动文件当前读写位置。
  \item Delete：删除文件，释放空间，删除目录项。
  \item Truncate：删除文件内容但保留属性。
  \item Open/Close：打开文件时把文件控制信息装入内存打开文件表，关闭时释放相关表项。
\end{points}
\plain{文件是用户看到的逻辑外存单位；用户不用管扇区，只管名字、属性和读写操作。}
\exam{常考文件属性、文件操作、open 为什么能减少目录搜索次数。}


\section{文件访问方法与逻辑结构}
\concept{顺序访问}
\begin{points}
  \item 顺序访问按文件中信息项的先后顺序读写。
  \item 适合文本流、日志、磁带等场景。
  \item 若要访问第 $n$ 个记录，通常需要从前面依次读到它。
\end{points}
\plain{文件怎么读有两层：逻辑上记录如何组织，物理上块如何放在磁盘。}
\exam{常考顺序/直接访问，顺序文件、索引文件、索引顺序文件、散列文件适用场景。}


\concept{直接访问/随机访问}
\begin{points}
  \item 直接访问允许按记录号或块号直接定位到某个位置。
  \item 磁盘支持随机访问，因此文件系统可实现直接访问。
  \item 适合数据库、索引文件、需要频繁定位的数据文件。
\end{points}
\plain{文件怎么读有两层：逻辑上记录如何组织，物理上块如何放在磁盘。}
\exam{常考顺序/直接访问，顺序文件、索引文件、索引顺序文件、散列文件适用场景。}


\concept{文件逻辑结构}
\begin{points}
  \item 无结构文件/流式文件：文件被看作字节序列，如普通文本和二进制文件。
  \item 有结构文件/记录式文件：文件由一组记录组成，每条记录有内部字段结构。
  \item 顺序文件：记录按某种顺序排列，适合顺序处理。
  \item 索引文件：建立索引表，从键值映射到记录位置，提高随机查找效率。
  \item 索引顺序文件：主文件按键有序，同时建立分组索引，兼顾顺序处理和快速定位。
  \item 散列文件：通过散列函数直接计算记录位置，适合等值查找。
\end{points}
\plain{文件怎么读有两层：逻辑上记录如何组织，物理上块如何放在磁盘。}
\exam{常考顺序/直接访问，顺序文件、索引文件、索引顺序文件、散列文件适用场景。}


\section{目录结构}
\concept{目录与目录项}
\begin{points}
  \item 目录用于组织系统中的文件并提供文件信息。
  \item 目录项通常包含文件名和文件控制块指针或 inode 号。
  \item 目录检索就是根据路径名逐级查找目录项，找到目标文件控制信息。
\end{points}
\plain{目录就是文件名到文件控制信息的映射表，路径就是沿目录树找到文件。}
\exam{常考一级/两级/树形目录优缺点、平均查找磁盘访问次数。}


\concept{一级目录}
\begin{points}
  \item 系统只有一个目录，所有文件都放在其中。
  \item 优点：实现简单。
  \item 缺点：文件多时查找慢，不能解决重名问题，不便于多用户管理。
\end{points}
\plain{目录就是文件名到文件控制信息的映射表，路径就是沿目录树找到文件。}
\exam{常考一级/两级/树形目录优缺点、平均查找磁盘访问次数。}


\concept{两级目录}
\begin{points}
  \item 每个用户有自己的用户文件目录 UFD，系统有主文件目录 MFD。
  \item 解决不同用户之间文件重名问题。
  \item 用户内部文件多时仍不够灵活。
\end{points}
\plain{目录就是文件名到文件控制信息的映射表，路径就是沿目录树找到文件。}
\exam{常考一级/两级/树形目录优缺点、平均查找磁盘访问次数。}


\concept{树形目录与路径}
\begin{points}
  \item 树形目录把目录组织成层次结构，便于分类管理。
  \item 绝对路径从根目录开始，如 \code{/home/a.txt}。
  \item 相对路径从当前目录开始。
  \item 优点：层次清晰、重名控制灵活、便于权限管理。
\end{points}
\plain{目录就是文件名到文件控制信息的映射表，路径就是沿目录树找到文件。}
\exam{常考一级/两级/树形目录优缺点、平均查找磁盘访问次数。}


\concept{无环图目录与共享}
\begin{points}
  \item 无环图目录允许多个目录项共享同一个文件或子目录，但不允许形成环。
  \item 共享文件需要引用计数，记录有多少目录项指向该文件。
  \item 引用计数为 0 时，文件数据和控制信息才能真正回收。
\end{points}
\plain{多个进程都想用同一批资源，所以必须规定能不能一起用、谁先用、用完怎么释放。}
\exam{常考互斥共享/同时共享举例，或者问并发和共享为什么是 OS 最基本特征。}


\section{文件系统安装、共享与保护}
\concept{文件系统安装 mount}
\begin{points}
  \item 文件系统必须挂载到某个目录位置后，用户才能通过统一目录树访问它。
  \item 挂载点是原有目录树中的一个目录。
  \item 挂载后，该目录下呈现被挂载文件系统的根内容。
\end{points}
\plain{文件是用户看到的逻辑外存单位；用户不用管扇区，只管名字、属性和读写操作。}
\exam{常考文件属性、文件操作、open 为什么能减少目录搜索次数。}


\concept{硬链接与符号链接}
\begin{points}
  \item 硬链接是多个目录项指向同一个 inode/文件控制块，本质上是同一个文件有多个名字。
  \item 删除一个硬链接只减少引用计数；只有引用计数为 0 时文件才真正删除。
  \item 符号链接是一个特殊文件，内容是目标文件路径。
  \item 符号链接可以跨文件系统，也可以指向目录；但目标删除后可能变成悬空链接。
\end{points}
\plain{硬链接共享同一个 inode，符号链接是另一个文件，内容是目标路径。打开文件表保存运行时读写状态。}
\exam{常考硬/软链接区别、引用计数、系统级/进程级打开文件表。}


\concept{文件保护}
\begin{points}
  \item 文件保护用于控制用户对文件的读、写、执行、删除等权限。
  \item 常见方法包括访问控制表 ACL、权限位、口令保护、能力表等。
  \item Unix 风格权限常分为所有者、同组用户、其他用户三类，每类有读写执行位。
\end{points}
\plain{把这个知识点放回本章主线理解：它回答的是 OS 为了管理资源、隐藏硬件、保证并发正确性而采用的某个对象、规则或步骤。}
\exam{常考选择/填空/简答，让你判断概念归属、比较相近概念，或按“定义-作用-机制-易错点”展开。}


\concept{打开文件表}
\begin{points}
  \item 打开文件后，系统把文件控制信息缓存在内存打开文件表中，避免每次读写都重新搜索目录。
  \item 系统级打开文件表记录文件位置、引用计数、磁盘位置等全局信息。
  \item 进程级打开文件表记录进程自己的文件描述符、访问模式、当前读写指针等。
  \item 若父子进程共享同一个打开文件表项，可能共享文件偏移；若分别打开同一文件，则偏移独立。
\end{points}
\plain{硬链接共享同一个 inode，符号链接是另一个文件，内容是目标路径。打开文件表保存运行时读写状态。}
\exam{常考硬/软链接区别、引用计数、系统级/进程级打开文件表。}


